\documentclass[11pt,a4paper]{article}
\usepackage{../templates/igcsephysics}

\renewcommand{\examcode}{Pre-IG}
\renewcommand{\papername}{Pre-IG Theory}
\renewcommand{\examseason}{June 2026}
\renewcommand{\examtime}{}
\renewcommand{\exammarks}{40}
\renewcommand{\examfooter}{Pre-IG Final Mock}
\renewcommand{\assetpath}{assets}

\begin{document}

\frontpage

\mcquestion{A digital stop-clock measures time in minutes and seconds.}
\mcqtext{The stop-clock reads 00:50 when it is started (i.e. 00 minutes 50 seconds).}
\mcqtext{It reads 02:10 when it is stopped.}
\mcqtext{What is the shortest possible time that has elapsed between starting and stopping the stop-clock?}
\mcqchoices
  {1 minute 20 seconds}
  {2 minutes 00 seconds}
  {2 minutes 10 seconds}
  {3 minutes 00 seconds}

\mcquestion{A long-distance runner wishes to calculate her average speed for a race.}
\mcqtext{Which calculation should she use?}
\mcqchoices
  {$\text{average speed}=\dfrac{\text{total distance}}{\text{total time}}$}
  {average speed = total distance $\times$ total time}
  {$\text{average speed}=\dfrac{\text{total time}}{\text{total distance}}$}
  {average speed = total distance + total time}

\newpage

\mcquestion{An arrow travels horizontally in a straight line at constant speed.}
\mcqtext{In which direction does the weight act?}
\mcqcenter{\includegraphics[width=105mm,trim=260 25 315 165,clip]{\assetpath/source_questions/q03.png}}

\mcquestion{The diagrams show four solid blocks with their dimensions and masses.}
\mcqtext{Which block has the greatest density?}
\mcqcenter{\includegraphics[width=120mm,trim=130 35 115 155,clip]{\assetpath/source_questions/q04.png}}

\newpage

\mcquestion{A wooden bar is pivoted at its centre so that it can rotate freely. Two equal forces $F$ are applied to the bar.}
\mcqtext{In which diagram is the turning effect greatest?}
\mcqcenter{\includegraphics[width=130mm,trim=40 0 35 190,clip]{\assetpath/source_questions/q05.png}}

\newpage

\mcquestion{A spring is suspended from a stand. Loads are added and the extensions are measured.}
\mcqcenter{\includegraphics[width=70mm,trim=395 810 420 210,clip]{\assetpath/source_questions/q06.png}}
\mcqtext{Which graph shows the result of plotting extension against load?}
\mcqcenter{\includegraphics[width=105mm,trim=120 45 145 670,clip]{\assetpath/source_questions/q06.png}}

\mcquestion{The diagram shows a stone suspended on a string under the surface of a liquid. The stone experiences a pressure caused by the liquid.}
\mcqcenter{\includegraphics[width=70mm,trim=650 355 660 220,clip]{\assetpath/source_questions/q07.png}}
\mcqtext{What would increase the pressure on the stone?}
\mcqchoices
  {decreasing the surface area of the stone}
  {increasing the mass of the stone}
  {lowering the stone deeper into the liquid}
  {using a liquid with a lower density}

\newpage

\mcquestion{A voltmeter and an ammeter are used to measure the resistance of a resistor.}
\mcqtext{Which diagram shows the voltmeter and the ammeter correctly connected?}
\mcqcenter{\includegraphics[width=118mm,trim=85 20 85 160,clip]{\assetpath/source_questions/q08.png}}

\mcquestion{Which object is in equilibrium?}
\mcqcenter{\includegraphics[width=110mm,trim=85 20 100 85,clip]{\assetpath/source_questions/q09.png}}

\mcquestion{Which diagram shows the dispersion of white light by a glass prism?}
\mcqcenter{\includegraphics[width=118mm,trim=100 30 100 120,clip]{\assetpath/source_questions/q10.png}}

\newpage

\mcquestion{A student draws a ray diagram to show how a ray of light is reflected by a number of mirrors.}
\mcqtext{Which reflection has not been drawn correctly?}
\mcqcenter{\includegraphics[width=110mm,trim=130 35 120 180,clip]{\assetpath/source_questions/q11.png}}

\mcquestion{The diagram shows four labelled changes of state between solid, liquid and gas.}
\mcqcenter{\includegraphics[width=95mm,trim=300 250 300 95,clip]{\assetpath/source_questions/q12.png}}
\mcqtext{Which changes need an energy input?}
\mcqchoiceswide
  {P and Q}
  {Q and R}
  {R and S}
  {S and P}

\newpage

\mcquestion{A student is designing a lighting circuit for a dolls' house. He sets up two different circuits.}
\mcqtext{Each circuit contains a 12V power supply and three identical lamps.}
\mcqtext{Each lamp is designed to operate at normal brightness when connected individually to a 12V supply.}
\mcqcenter{\includegraphics[width=120mm,trim=120 395 100 330,clip]{\assetpath/source_questions/q13.png}}
\mcqtext{Which statement is correct?}
\mcqchoices
  {In circuit 1, each of the lamps is at normal brightness.}
  {In circuit 1, if one lamp fails, the other lamps remain lit.}
  {In circuit 2, if one lamp fails, the other lamps remain lit.}
  {In circuit 2, the current from the power supply is less than in circuit 1.}

\newpage

\questionpart{a}{The arrows on Fig. \figref{fig:changes-state} represent changes of state.}{3}
\examfigure{fig:changes-state}{%
  \includegraphics[width=135mm,trim=40 205 80 85,clip]{\assetpath/source_questions/q14.png}%
}
\parttext{One of the arrows is labelled. Label each of the other arrows with the correct change of state.}
\parttext{Write the change of state on the dotted lines next to each arrow.}
\marksright{3}
\totalmarks

\newpage

\questionpart{a}{A student has a metal object.}{10}

\partsubquestion{a}{i}{The student measures the mass of the object.}
\subparttext{State the name of the equipment used to measure the mass.}
\subpartwrittenanswer{ }{1}

\subpartquestion{ii}{The mass of the metal object is 1260g. The volume of the metal is 150cm\textsuperscript{3}.}
\subparttext{Calculate the density of the metal. Include the unit.}
\working[34mm]
\answer{density}{}{4}

\subpartquestion{iii}{The mass of the metal object is given in grams. State the mass in kg.}
\working[18mm]
\answer{mass}{kg}{1}

\partquestion{b}{A vase is placed on a table. Forces X and Y act on the vase, as shown in Fig. \figref{fig:vase-forces}.}
\examfigure{fig:vase-forces}{%
  \includegraphics[width=68mm,trim=470 900 480 980,clip]{\assetpath/source_questions/q15.png}%
}
\parttext{The mass of the vase is 0.25kg. The vase is not moving.}
\parttext{Calculate the value of force X and the value of force Y.}
\working[22mm]
\rightanswerline[70mm]{X}
\rightanswerline[70mm]{Y}
\rightanswermark{4}
\totalmarks

\newpage

\questionpart{a}{A student rubs a plastic rod with a dry cloth, as shown in Fig. \figref{fig:plastic-rod}. The rod becomes negatively charged.}{6}
\examfigure{fig:plastic-rod}{%
  \includegraphics[width=95mm,trim=330 2470 380 165,clip]{\assetpath/source_questions/q16_q17.png}%
}

\partsubquestion{a}{i}{Use words from the box to complete the sentence.}
\begin{center}
\renewcommand{\arraystretch}{1.7}
\begin{tabular}{|c@{\hspace{9mm}}c@{\hspace{9mm}}c@{\hspace{9mm}}c@{\hspace{9mm}}c@{\hspace{9mm}}c|}
\hline
air & cloth & electrons & hand & neutrons & protons \\ \hline
\end{tabular}
\end{center}
\subparttext{The rod becomes negatively charged because \answerline[42mm] move from the}
\subpartwrittenanswer[92mm]{to the rod.}{2}

\subpartquestion{ii}{The student moves the rod close to a suspended, charged rod. The two rods repel each other.}
\subparttext{State the type of charge on the suspended rod.}
\subpartwrittenanswer{ }{1}

\subpartquestion{iii}{Explain your answer to (a)(ii).}
\subpartwrittenline{ }
\subpartwrittenanswer{ }{1}

\partquestion{b}{A device has a metal case. Any charge on the case must be able to move to earth.}
\partsubquestion{b}{i}{Draw one ring around a material that is suitable for the connection to earth.}
\begin{center}
\begin{tabular}{c@{\hspace{18mm}}c@{\hspace{18mm}}c@{\hspace{18mm}}c}
copper & glass & plastic & rubber
\end{tabular}
\end{center}
\marksright{1}

\subpartquestion{ii}{Explain your answer to (b)(i).}
\subpartwrittenline{ }
\subpartwrittenanswer{ }{1}
\totalmarks

\newpage

\question{Fig. \figref{fig:mirror-reflection} shows a ray of light that is reflected by a mirror.}{6}
\examfigure{fig:mirror-reflection}{%
  \includegraphics[width=95mm,trim=300 2080 300 100,clip]{\assetpath/source_questions/q18.png}%
}

\partsubquestion{a}{i}{State the name of line X shown on Fig. \figref{fig:mirror-reflection}.}
\subpartwrittenanswer{ }{1}

\subpartquestion{ii}{State the name of angle Y shown on Fig. \figref{fig:mirror-reflection}.}
\subpartwrittenanswer{ }{1}

\subpartquestion{iii}{A student moves the ray of light and doubles the size of angle Y. State the effect on angle Z.}
\subpartwrittenanswer{ }{1}

\partquestion{b}{Fig. \figref{fig:lens-image} shows a converging lens used to form an image I of an object O.}
\examfigure{fig:lens-image}{%
  \includegraphics[width=130mm,trim=30 560 65 1230,clip]{\assetpath/source_questions/q18.png}%
}

\partsubquestion{b}{i}{State the name of the points labelled F on Fig. \figref{fig:lens-image}.}
\subpartwrittenanswer{ }{1}

\subpartquestion{ii}{Describe the nature of the image I.}
\subpartwrittenline{ }
\subpartwrittenline{ }
\subpartwrittenanswer{ }{2}
\totalmarks

\newpage

\questionpart{a}{State the two conditions which must be true for an object to be in equilibrium.}{10}
\rightanswerline[120mm]{condition 1}
\rightanswerline[120mm]{condition 2}
\rightanswermark{2}

\partquestion{b}{Fig. \figref{fig:metre-rule} shows a uniform metre rule PQ in equilibrium.}
\examfigure{fig:metre-rule}{%
  \includegraphics[width=115mm,trim=35 742 35 250,clip]{\assetpath/source_questions/q19.png}%
}
\parttext{The distance PQ is 100cm. The mass of the metre rule is 0.12kg and its weight is W.}

\partsubquestion{b}{i}{On Fig. \figref{fig:metre-rule}, draw and label:}
\subparttext{1. an arrow to show the force W acting on PQ at the centre of mass}
\subparttext{2. an arrow to show the force R acting on PQ at the pivot.}
\marksright{2}

\subpartquestion{ii}{By taking moments about the pivot, calculate F.}
\working[28mm]
\answer{F}{}{4}

\subpartquestion{iii}{Calculate R.}
\working[22mm]
\answer{R}{}{2}
\totalmarks

\end{document}
